\documentclass[11pt]{article}
\usepackage[margin=1in]{geometry}
\usepackage{amsmath, amssymb, amsthm}
\usepackage{booktabs}
\usepackage{graphicx}

\newtheorem{theorem}{Theorem}
\newtheorem{lemma}{Lemma}

\title{Project Euler Problem 394: Eating Pie}
\author{}
\date{}

\begin{document}
\maketitle

\section{Problem Statement}

Jeff eats a circular pie by making an initial radial cut, then repeatedly:
\begin{enumerate}
    \item Choosing one of the two exposed edges uniformly at random,
    \item Cutting at angle $\alpha \sim \text{Uniform}(0, \theta)$ from that edge (where $\theta$ is the remaining angle),
    \item Eating the resulting slice.
\end{enumerate}

Let $E(n)$ be the expected number of slices Jeff eats until the remaining pie has angle less than $1/n$ of the original. Find:
\[
    \lceil E(10) \rceil + \lceil E(100) \rceil + \lceil E(1000) \rceil.
\]

\section{Solution}

\subsection{Reduction to Uniform Products}

Normalize the full pie angle to 1. After each cut, the remaining fraction satisfies:
\[
    \theta_{k+1} = \theta_k - \alpha_k, \qquad \alpha_k \sim \text{Uniform}(0, \theta_k).
\]

By symmetry, the choice of edge does not affect the distribution. Writing $V_k = 1 - \alpha_k/\theta_k \sim \text{Uniform}(0,1)$:
\[
    \theta_k = \prod_{i=1}^{k} V_i, \qquad V_i \overset{\text{iid}}{\sim} \text{Uniform}(0,1).
\]

\subsection{Reduction to Exponential Sums}

We need $N = \min\{k : \theta_k < 1/n\}$, i.e.,
\[
    N = \min\left\{k : \sum_{i=1}^{k} (-\ln V_i) > \ln n\right\}.
\]

Since $X_i = -\ln V_i \sim \text{Exp}(1)$, this is the first passage time of a standard Poisson process past level $t = \ln n$.

\subsection{Exact Expectation}

\begin{theorem}
    If $X_1, X_2, \ldots \overset{\text{iid}}{\sim} \text{Exp}(1)$ and $S_k = \sum_{i=1}^k X_i$, then for $N = \min\{k : S_k > t\}$ with $t > 0$:
    \[
        E[N] = 1 + t.
    \]
\end{theorem}

\begin{proof}
    \begin{align*}
        E[N] &= \sum_{k=0}^{\infty} P(N > k) = \sum_{k=0}^{\infty} P(S_k \leq t).
    \end{align*}
    Since $S_k \sim \text{Gamma}(k,1)$ and using the Poisson--Gamma duality:
    \[
        P(S_k \leq t) = P\bigl(\text{Poisson}(t) \geq k\bigr) = \sum_{j=k}^{\infty} \frac{e^{-t} t^j}{j!}.
    \]
    Therefore:
    \begin{align*}
        E[N] &= \sum_{k=0}^{\infty} \sum_{j=k}^{\infty} \frac{e^{-t} t^j}{j!}
             = \sum_{j=0}^{\infty} (j+1) \frac{e^{-t} t^j}{j!} \\
             &= e^{-t} \left( \sum_{j=0}^{\infty} \frac{j \, t^j}{j!} + \sum_{j=0}^{\infty} \frac{t^j}{j!} \right)
             = e^{-t}(t e^t + e^t)
             = t + 1.
    \end{align*}
\end{proof}

\subsection{Final Computation}

With $E(n) = 1 + \ln n$:

\begin{center}
\begin{tabular}{@{}cccc@{}}
    \toprule
    $n$ & $E(n) = 1 + \ln n$ & $\lceil E(n) \rceil$ \\
    \midrule
    10   & 3.30259 & 4 \\
    100  & 5.60517 & 6 \\
    1000 & 7.90776 & 8 \\
    \bottomrule
\end{tabular}
\end{center}

\section{Answer}

\[
\boxed{3.2370342194}
\]

\end{document}
